\chapter{Black--Scholes from the binomial tree, every step}\label{app:bs}

% Two columns: the line of mathematics on the left, how it follows from the line before on the right.
\newenvironment{deriv}{\par\smallskip\noindent\renewcommand{\arraystretch}{1.3}%
  \begin{tabular}{@{}p{0.58\linewidth}@{\hspace{0.03\linewidth}}p{0.39\linewidth}@{}}}{\end{tabular}\par\smallskip}
\newcommand{\dl}[2]{\mbox{$\displaystyle #1$} & \leavevmode\raggedright\small\color{thmc!80!black}#2\tabularnewline[1.2ex]}

This appendix is the full version of Problem Q2(a) of the April 2024 paper. It starts with one period and a portfolio that copies the option, and it ends with the Black--Scholes formula. Part E then does the put by the same steps, and Part F builds the copying portfolio at every node of the tree and shows that it is self-financing. Each row does one thing. The right-hand column says how that row follows from the row before it.

\begin{setup}[The market]
A stock starts at $S_0$. Each period it is multiplied by $u$ (up) or by $d$ (down), with $d<u$. Money in the bond grows by $1+r$ each period. A European call with strike $K$ pays $(S_T-K)^+=\max(S_T-K,0)$ at time $T$.
\end{setup}

\section*{Part A: one period}

\begin{steps}
\item A claim pays $X_u$ if the stock goes up and $X_d$ if it goes down. Buy $\Delta$ shares and put $b$ in the bond, and choose $\Delta$ and $b$ so that the portfolio pays exactly the claim in both cases:
\begin{deriv}
\dl{\Delta\,uS_0+b(1+r)=X_u}{up: each share is worth $uS_0$, the bond has grown to $b(1+r)$}
\dl{\Delta\,dS_0+b(1+r)=X_d}{down: each share is worth $dS_0$}
\end{deriv}

\item Find $\Delta$:
\begin{deriv}
\dl{\Delta(u-d)S_0=X_u-X_d}{subtract the second equation from the first; $b(1+r)$ cancels}
\dl{\Delta=\frac{X_u-X_d}{(u-d)S_0}}{divide by $(u-d)S_0$}
\end{deriv}

\item Find $b$:
\begin{deriv}
\dl{b(1+r)=X_u-u\,\Delta S_0}{the up equation of Step~1, solved for $b(1+r)$}
\dl{\phantom{b(1+r)}=X_u-\frac{u(X_u-X_d)}{u-d}}{put in $\Delta S_0=\frac{X_u-X_d}{u-d}$ from Step~2}
\dl{\phantom{b(1+r)}=\frac{(u-d)X_u-uX_u+uX_d}{u-d}}{common denominator $u-d$}
\dl{\phantom{b(1+r)}=\frac{uX_d-dX_u}{u-d}}{$(u-d)X_u-uX_u=-dX_u$}
\dl{b=\frac{uX_d-dX_u}{(1+r)(u-d)}}{divide by $1+r$}
\end{deriv}

\item The portfolio pays the same as the claim in both cases, so it must cost the same today; otherwise buying the cheaper one and selling the dearer one is an arbitrage. So the claim's price is the portfolio's cost:
\begin{deriv}
\dl{X^*=\Delta S_0+b}{$\Delta$ shares at $S_0$ each, plus $b$ in the bond}
\dl{\phantom{X^*}=\frac{X_u-X_d}{u-d}+\frac{uX_d-dX_u}{(1+r)(u-d)}}{Steps 2 and 3}
\dl{\phantom{X^*}=\frac{(1+r)(X_u-X_d)+uX_d-dX_u}{(1+r)(u-d)}}{common denominator $(1+r)(u-d)$}
\dl{\phantom{X^*}=\frac{(1+r-d)X_u+\big(u-(1+r)\big)X_d}{(1+r)(u-d)}}{collect the $X_u$ terms and the $X_d$ terms}
\dl{\phantom{X^*}=\frac{p^*X_u+(1-p^*)X_d}{1+r}}{name $p^*=\frac{1+r-d}{u-d}$; then $1-p^*=\frac{u-(1+r)}{u-d}$}
\end{deriv}
This is \cardrow{card:finite}{1}. The price is an average of the two payoffs with weights $p^*$ and $1-p^*$, divided by $1+r$.

\item Two facts about $p^*$ that we use later:
\begin{deriv}
\dl{0<p^*<1\iff d<1+r<u}{$p^*>0$ needs $1+r>d$; $1-p^*>0$ needs $u>1+r$}
\dl{up^*+d(1-p^*)=d+p^*(u-d)}{multiply out: $up^*+d-dp^*$}
\dl{\phantom{up^*+d(1-p^*)}=d+(1+r-d)=1+r\qquad(\star)}{$p^*(u-d)=1+r-d$ by the definition of $p^*$}
\end{deriv}
So $(\star)$ says that under $p^*$ the stock grows on average at the bond rate.
\end{steps}

\section*{Part B: $T$ periods}

\begin{steps}[resume]
\item Every node of the tree is a one-period market, so Step~4 applies at each node. Let $V_t[s]$ be the claim's value at time $t$ when $S_t=s$:
\begin{deriv}
\dl{V_T[s]=g(s)}{at the end the claim is worth its payoff $g$}
\dl{V_t[s]=\frac{p^*V_{t+1}[us]+(1-p^*)V_{t+1}[ds]}{1+r}}{Step~4 with $X_u=V_{t+1}[us]$, $X_d=V_{t+1}[ds]$ (\cref{card:tree}, box)}
\end{deriv}

\item Two periods first, to see the pattern:
\begin{deriv}
\dl{V_1[uS_0]=\frac{p^*g(u^2S_0)+(1-p^*)g(udS_0)}{1+r}}{Step~6 at the up node}
\dl{V_1[dS_0]=\frac{p^*g(udS_0)+(1-p^*)g(d^2S_0)}{1+r}}{Step~6 at the down node}
\dl{V_0=\frac{p^*V_1[uS_0]+(1-p^*)V_1[dS_0]}{1+r}}{Step~6 at time $0$}
\dl{\phantom{V_0}=\frac1{(1+r)^2}\Big[(p^*)^2g(u^2S_0)}{put in the two lines above and collect}
\dl{\qquad+2p^*(1-p^*)g(udS_0)+(1-p^*)^2g(d^2S_0)\Big]}{}
\end{deriv}
The $2$ counts the two paths (up then down, down then up) that both end at $udS_0$.

\item $T$ periods. Each step back multiplies a path by $p^*$ for an up-move or $1-p^*$ for a down-move, and divides by $1+r$. After $T$ steps back:
\begin{deriv}
\dl{V_0=\frac1{(1+r)^T}\sum_{\text{paths}}(\text{weight of the path})}{Step~7's pattern, repeated $T$ times}
\dl{\qquad\times g(\text{end price of the path})}{}
\dl{\text{a path with $j$ ups: weight }(p^*)^j(1-p^*)^{T-j}}{$j$ factors $p^*$ and $T-j$ factors $1-p^*$}
\dl{\text{it ends at }u^jd^{T-j}S_0}{the order of the moves does not change the product}
\dl{\text{there are }\binom Tj\text{ such paths}}{choose which $j$ of the $T$ periods are ups}
\dl{V_0=\frac1{(1+r)^T}\sum_{j=0}^T\binom Tj(p^*)^j(1-p^*)^{T-j}}{group the paths by $j$ (\cardrow{card:tree}{1})}
\dl{\qquad\times g\big(u^jd^{T-j}S_0\big)}{}
\end{deriv}

\item Now the call, $g(s)=(s-K)^+$. The end price $u^jd^{T-j}S_0$ grows with $j$, because each extra up replaces a $d$ by a $u>d$. Let $m$ be the least $j$ with $u^jd^{T-j}S_0>K$:
\begin{deriv}
\dl{j<m:\ (u^jd^{T-j}S_0-K)^+=0}{the end price is at most $K$}
\dl{j\ge m:\ (u^jd^{T-j}S_0-K)^+=u^jd^{T-j}S_0-K}{the end price is above $K$}
\dl{C=\frac1{(1+r)^T}\sum_{j\ge m}\binom Tj(p^*)^j(1-p^*)^{T-j}}{Step~8 with the $j<m$ terms dropped}
\dl{\qquad\times\big(u^jd^{T-j}S_0-K\big)}{}
\end{deriv}

\item Split the bracket into the stock part and the $K$ part:
\begin{deriv}
\dl{C=S_0\sum_{j\ge m}\binom Tj\frac{(p^*)^j(1-p^*)^{T-j}u^jd^{T-j}}{(1+r)^T}}{the $u^jd^{T-j}S_0$ part; $S_0$ comes out of the sum}
\dl{\qquad-\frac{K}{(1+r)^T}\sum_{j\ge m}\binom Tj(p^*)^j(1-p^*)^{T-j}}{the $K$ part; $K$ and $(1+r)^{-T}$ come out}
\end{deriv}

\item Tidy the stock part:
\begin{deriv}
\dl{\frac{(p^*)^j(1-p^*)^{T-j}u^jd^{T-j}}{(1+r)^T}}{the term inside the first sum of Step~10}
\dl{\qquad=\Big(\frac{up^*}{1+r}\Big)^j\Big(\frac{d(1-p^*)}{1+r}\Big)^{T-j}}{write $(1+r)^T=(1+r)^j(1+r)^{T-j}$ and pair each factor}
\dl{\hat p=\frac{up^*}{1+r}}{a name for the first bracket}
\dl{1-\hat p=\frac{1+r-up^*}{1+r}=\frac{d(1-p^*)}{1+r}}{$1+r-up^*=d(1-p^*)$ by $(\star)$ of Step~5}
\dl{\frac{(p^*)^j(1-p^*)^{T-j}u^jd^{T-j}}{(1+r)^T}=\hat p^{\,j}(1-\hat p)^{T-j}}{put both brackets back}
\end{deriv}
And $0<\hat p<1$: by $(\star)$, $up^*$ and $d(1-p^*)$ are two positive numbers that add up to $1+r$. So $\hat p$ is a probability too.

\item The binomial call price (\cardrow{card:tree}{3}):
\begin{deriv}
\dl{C=S_0\sum_{j\ge m}\binom Tj\hat p^{\,j}(1-\hat p)^{T-j}}{Step~10, with Step~11 inside the first sum}
\dl{\qquad-\frac{K}{(1+r)^T}\sum_{j\ge m}\binom Tj(p^*)^j(1-p^*)^{T-j}}{the $K$ part of Step~10, unchanged}
\end{deriv}
\end{steps}

\section*{Part C: read the two sums as probabilities}

\begin{steps}[resume]
\item Let $J$ be the number of up-moves. When each period is up with probability $p$, independently, $J\sim\mathrm{Bin}(T,p)$, so $P(J\ge m)=\sum_{j\ge m}\binom Tjp^j(1-p)^{T-j}$:
\begin{deriv}
\dl{\sum_{j\ge m}\binom Tj(p^*)^j(1-p^*)^{T-j}=P^*(J\ge m)}{$P^*$: each period is up with probability $p^*$}
\dl{\sum_{j\ge m}\binom Tj\hat p^{\,j}(1-\hat p)^{T-j}=\hat P(J\ge m)}{$\hat P$: each period is up with probability $\hat p$}
\dl{J\ge m\iff S_T>K}{$S_T=u^Jd^{T-J}S_0$ and the definition of $m$ in Step~9}
\dl{C=S_0\,\hat P(S_T>K)-\frac{K}{(1+r)^T}\,P^*(S_T>K)}{put the three lines into Step~12}
\end{deriv}
Part D lets the periods shrink and finds each of the two probabilities.
\end{steps}

\section*{Part D: let the periods shrink}

\begin{steps}[resume]
\item Fix the real time to expiry $\tau$. Cut it into $n$ periods of length $\Delta t=\tau/n$, and choose the tree to match a stock with volatility $\sigma$ and a bank paying the continuous rate $r$:
\begin{deriv}
\dl{u=e^{\sigma\sqrt{\Delta t}},\qquad d=e^{-\sigma\sqrt{\Delta t}}}{up and down moves of size $\sigma\sqrt{\Delta t}$ in the log-price}
\dl{1+r_n=e^{r\Delta t}}{one period of continuous interest at rate $r$}
\dl{(1+r_n)^n=e^{rn\Delta t}=e^{r\tau}}{$n\Delta t=\tau$}
\dl{\frac{K}{(1+r_n)^n}=Ke^{-r\tau}}{the discount in Step~13 with $T=n$}
\end{deriv}
The condition $d<1+r_n<u$ of Step~5 is $-\sigma\sqrt{\Delta t}<r\Delta t<\sigma\sqrt{\Delta t}$, which holds once $\Delta t$ is small.

\item The log-price is a random walk:
\begin{deriv}
\dl{\ln\frac{S_\tau}{S_0}=J\,\sigma\sqrt{\Delta t}-(n-J)\,\sigma\sqrt{\Delta t}}{$S_\tau=u^Jd^{n-J}S_0$; take logs}
\end{deriv}
So every $\Delta t$ the log-price moves up or down by $\Delta x=\sigma\sqrt{\Delta t}$. This is the walk of \cref{card:scaledwalk}.

\item Taylor expansions, using $e^x=1+x+\frac12x^2+O(x^3)$:
\begin{deriv}
\dl{u=1+\sigma\sqrt{\Delta t}+\tfrac12\sigma^2\Delta t+O(\Delta t^{3/2})}{$x=\sigma\sqrt{\Delta t}$, so $x^2=\sigma^2\Delta t$}
\dl{d=1-\sigma\sqrt{\Delta t}+\tfrac12\sigma^2\Delta t+O(\Delta t^{3/2})}{$x=-\sigma\sqrt{\Delta t}$; the square is the same}
\dl{e^{r\Delta t}=1+r\Delta t+O(\Delta t^2)}{$x=r\Delta t$, so $x^2$ is of order $\Delta t^2$}
\end{deriv}

\item Expand $p^*$:
\begin{deriv}
\dl{p^*=\frac{e^{r\Delta t}-d}{u-d}}{Step~4 with $1+r_n=e^{r\Delta t}$}
\dl{e^{r\Delta t}-d=\sigma\sqrt{\Delta t}+(r-\tfrac12\sigma^2)\Delta t+O(\Delta t^{3/2})}{Step~16: subtract; the $1$s cancel}
\dl{u-d=2\sigma\sqrt{\Delta t}+O(\Delta t^{3/2})}{Step~16: subtract; the $\frac12\sigma^2\Delta t$ terms cancel}
\dl{p^*=\frac{\sigma\sqrt{\Delta t}+(r-\frac12\sigma^2)\Delta t}{2\sigma\sqrt{\Delta t}}+O(\Delta t)}{divide the two lines}
\dl{\phantom{p^*}=\frac12\Big(1+\frac{r-\frac12\sigma^2}{\sigma}\sqrt{\Delta t}\Big)+O(\Delta t)}{split the fraction and cancel $\sigma\sqrt{\Delta t}$}
\end{deriv}

\item Compare with \cref{card:scaledwalk}: there $p=\frac12\big(1+\frac\mu\sigma\sqrt{\Delta t}\big)$. So under $P^*$ the walk has drift $\mu=r-\frac12\sigma^2$. Check the mean and the variance with the card's box:
\begin{deriv}
\dl{2p^*-1=\frac\mu\sigma\sqrt{\Delta t}+O(\Delta t)}{Step~17: double it and subtract $1$}
\dl{\text{mean}=\frac{\tau}{\Delta t}\,\sigma\sqrt{\Delta t}\,(2p^*-1)}{\cref{card:scaledwalk}, box, with $t=\tau$, $\Delta x=\sigma\sqrt{\Delta t}$}
\dl{\phantom{\text{mean}}=\mu\tau+O(\sqrt{\Delta t})\ \to\ \mu\tau}{$\frac{\tau}{\Delta t}\sigma\sqrt{\Delta t}\cdot\frac\mu\sigma\sqrt{\Delta t}=\mu\tau$; the $O(\Delta t)$ leftover gives $\frac{\tau}{\Delta t}\sigma\sqrt{\Delta t}\,O(\Delta t)=O(\sqrt{\Delta t})$}
\dl{\text{variance}=\frac{\tau}{\Delta t}\,\sigma^2\Delta t\big(1-(2p^*-1)^2\big)}{\cref{card:scaledwalk}, box}
\dl{\phantom{\text{variance}}=\sigma^2\tau\big(1-O(\Delta t)\big)\ \to\ \sigma^2\tau}{$(2p^*-1)^2=O(\Delta t)$ by the first line}
\end{deriv}

\item The limit is normal. $\ln(S_\tau/S_0)$ is a sum of $n$ independent steps, so by the central limit theorem (taken as known, as in \cref{card:scaledwalk}) it becomes normal with the mean and variance of Step~18:
\begin{deriv}
\dl{\text{under }P^*:\quad\ln\frac{S_\tau}{S_0}\ \to\ N\big((r-\tfrac12\sigma^2)\tau,\ \sigma^2\tau\big)}{Step~18 with $\mu=r-\frac12\sigma^2$}
\end{deriv}

\item The $K$ part. Use \cref{card:normal-prob}: if $\xi\sim N(a,s^2)$ then $P(\xi>c)=1-\Phi\big(\frac{c-a}{s}\big)=\Phi\big(\frac{a-c}{s}\big)$:
\begin{deriv}
\dl{P^*(S_\tau>K)=P^*\Big(\ln\frac{S_\tau}{S_0}>\ln\frac K{S_0}\Big)}{take logs of both sides; $\ln$ is increasing}
\dl{\phantom{P^*(S_\tau>K)}\to\Phi\bigg(\frac{(r-\frac12\sigma^2)\tau-\ln(K/S_0)}{\sigma\sqrt\tau}\bigg)}{Step~19 and the card, with $a=(r-\frac12\sigma^2)\tau$, $s=\sigma\sqrt\tau$, $c=\ln\frac K{S_0}$}
\dl{\phantom{P^*(S_\tau>K)}=\Phi\bigg(\frac{\ln(S_0/K)+(r-\frac12\sigma^2)\tau}{\sigma\sqrt\tau}\bigg)}{$-\ln\frac K{S_0}=\ln\frac{S_0}K$}
\dl{\phantom{P^*(S_\tau>K)}=\Phi(h-\sigma\sqrt\tau)}{name $h=\frac{\ln(S_0/K)+(r+\frac12\sigma^2)\tau}{\sigma\sqrt\tau}$; subtracting $\sigma\sqrt\tau=\frac{\sigma^2\tau}{\sigma\sqrt\tau}$ turns $+\frac12\sigma^2$ into $-\frac12\sigma^2$}
\end{deriv}
That the probability of the event $J\ge m$ converges to the normal probability is also taken as known; it holds because the normal limit has no jumps.

\item Expand $\hat p$ the same way:
\begin{deriv}
\dl{\hat p=\frac{up^*}{1+r_n}=e^{\sigma\sqrt{\Delta t}-r\Delta t}\,p^*}{Step~11 with $u=e^{\sigma\sqrt{\Delta t}}$ and $1+r_n=e^{r\Delta t}$}
\dl{e^{\sigma\sqrt{\Delta t}-r\Delta t}=1+\sigma\sqrt{\Delta t}+O(\Delta t)}{$e^x=1+x+O(x^2)$ with $x=\sigma\sqrt{\Delta t}-r\Delta t$}
\dl{\hat p=\big(1+\sigma\sqrt{\Delta t}\big)\cdot\frac12\Big(1+\frac{r-\frac12\sigma^2}{\sigma}\sqrt{\Delta t}\Big)+O(\Delta t)}{Step~17 for $p^*$}
\dl{\phantom{\hat p}=\frac12\Big(1+\sigma\sqrt{\Delta t}+\frac{r-\frac12\sigma^2}{\sigma}\sqrt{\Delta t}\Big)+O(\Delta t)}{multiply out; the product of the two $\sqrt{\Delta t}$ terms is $O(\Delta t)$}
\dl{\phantom{\hat p}=\frac12\Big(1+\frac{r+\frac12\sigma^2}{\sigma}\sqrt{\Delta t}\Big)+O(\Delta t)}{$\sigma=\frac{\sigma^2}{\sigma}$, and $\sigma^2+r-\frac12\sigma^2=r+\frac12\sigma^2$}
\end{deriv}

\item This is the walk of \cref{card:scaledwalk} with drift $\hat\mu=r+\frac12\sigma^2$. Steps 18 and 19 go through word for word with $\hat\mu$ in place of $\mu$:
\begin{deriv}
\dl{\text{under }\hat P:\quad\ln\frac{S_\tau}{S_0}\ \to\ N\big((r+\tfrac12\sigma^2)\tau,\ \sigma^2\tau\big)}{Steps 18--19 with $\mu=r+\frac12\sigma^2$}
\end{deriv}

\item The stock part, as in Step~20:
\begin{deriv}
\dl{\hat P(S_\tau>K)=\hat P\Big(\ln\frac{S_\tau}{S_0}>\ln\frac K{S_0}\Big)}{take logs}
\dl{\phantom{\hat P(S_\tau>K)}\to\Phi\bigg(\frac{\ln(S_0/K)+(r+\frac12\sigma^2)\tau}{\sigma\sqrt\tau}\bigg)}{Step~22 and \cref{card:normal-prob}, with $a=(r+\frac12\sigma^2)\tau$}
\dl{\phantom{\hat P(S_\tau>K)}=\Phi(h)}{the $h$ of Step~20}
\end{deriv}

\item Put the parts together:
\begin{deriv}
\dl{C=S_0\,\hat P(S_\tau>K)-\frac{K}{(1+r_n)^n}\,P^*(S_\tau>K)}{Step~13 with $T=n$}
\dl{\phantom{C}\to S_0\,\Phi(h)-Ke^{-r\tau}\,\Phi(h-\sigma\sqrt\tau)}{Steps 23, 14 and 20}
\end{deriv}
\end{steps}

\answer{$C=S_0\Phi(h)-Ke^{-r\tau}\Phi(h-\sigma\sqrt\tau)$ with $h=\dfrac{\ln(S_0/K)+(r+\frac12\sigma^2)\tau}{\sigma\sqrt\tau}$. This is the Black--Scholes formula (\cardrow{card:bsprice}{1}). The limit of the tree is \cardrow{card:bsprice}{8}.}

\section*{Part E: the put, by the same steps}

A European put pays $(K-S_T)^+=\max(K-S_T,0)$. Steps 1--8 never used the payoff, so they hold as they are. Only Steps 9--13 change, and then Steps 20 and 23 are used through their complements.

\begin{steps}[resume]
\item Step~9 for the put, $g(s)=(K-s)^+$. Now the paths that pay are the ones that end at or below $K$:
\begin{deriv}
\dl{j\ge m:\ (K-u^jd^{T-j}S_0)^+=0}{the end price is above $K$ (definition of $m$, Step~9)}
\dl{j<m:\ (K-u^jd^{T-j}S_0)^+=K-u^jd^{T-j}S_0}{the end price is at most $K$}
\dl{P=\frac1{(1+r)^T}\sum_{j<m}\binom Tj(p^*)^j(1-p^*)^{T-j}}{Step~8 with the $j\ge m$ terms dropped}
\dl{\qquad\times\big(K-u^jd^{T-j}S_0\big)}{}
\end{deriv}

\item Split and tidy, as in Steps 10 and 11:
\begin{deriv}
\dl{P=\frac{K}{(1+r)^T}\sum_{j<m}\binom Tj(p^*)^j(1-p^*)^{T-j}}{the $K$ part; $K$ and $(1+r)^{-T}$ come out}
\dl{\qquad-S_0\sum_{j<m}\binom Tj\hat p^{\,j}(1-\hat p)^{T-j}}{the stock part; the term inside is the one tidied in Step~11}
\end{deriv}

\item Read the sums as probabilities, as in Step~13:
\begin{deriv}
\dl{\sum_{j<m}\binom Tj(p^*)^j(1-p^*)^{T-j}=P^*(J<m)}{$J\sim\mathrm{Bin}(T,p^*)$ under $P^*$}
\dl{\sum_{j<m}\binom Tj\hat p^{\,j}(1-\hat p)^{T-j}=\hat P(J<m)}{$J\sim\mathrm{Bin}(T,\hat p)$ under $\hat P$}
\dl{J<m\iff S_T\le K}{the opposite of $J\ge m\iff S_T>K$ (Step~13)}
\dl{P=\frac{K}{(1+r)^T}\,P^*(S_T\le K)-S_0\,\hat P(S_T\le K)}{put the three lines into Step~26}
\end{deriv}

\item Let the periods shrink. $\{S_\tau\le K\}$ is the opposite of $\{S_\tau>K\}$, so each probability is $1$ minus the call's, and \cref{card:normal-prob} gives $1-\Phi(x)=\Phi(-x)$:
\begin{deriv}
\dl{P^*(S_\tau\le K)=1-P^*(S_\tau>K)}{complement}
\dl{\phantom{P^*(S_\tau\le K)}\to1-\Phi(h-\sigma\sqrt\tau)}{Step~20}
\dl{\phantom{P^*(S_\tau\le K)}=\Phi(\sigma\sqrt\tau-h)}{$1-\Phi(x)=\Phi(-x)$ with $x=h-\sigma\sqrt\tau$}
\dl{\hat P(S_\tau\le K)=1-\hat P(S_\tau>K)}{complement}
\dl{\phantom{\hat P(S_\tau\le K)}\to1-\Phi(h)=\Phi(-h)}{Step~23, then the same rule}
\end{deriv}

\item Put the parts together:
\begin{deriv}
\dl{P=\frac{K}{(1+r_n)^n}\,P^*(S_\tau\le K)-S_0\,\hat P(S_\tau\le K)}{Step~27 with $T=n$}
\dl{\phantom{P}\to Ke^{-r\tau}\,\Phi(\sigma\sqrt\tau-h)-S_0\,\Phi(-h)}{Steps 14 and 28}
\end{deriv}

\item Check with put--call parity (\cref{card:parity}, box): $S_0+P-C=Ke^{-r\tau}$, i.e.\ $C-P=S_0-Ke^{-r\tau}$:
\begin{deriv}
\dl{C-P=S_0\big[\Phi(h)+\Phi(-h)\big]}{Step~24 minus Step~29, the $S_0$ terms}
\dl{\qquad-Ke^{-r\tau}\big[\Phi(h-\sigma\sqrt\tau)+\Phi(\sigma\sqrt\tau-h)\big]}{the $K$ terms}
\dl{\phantom{C-P}=S_0-Ke^{-r\tau}}{$\Phi(x)+\Phi(-x)=1$ (\cref{card:normal-prob}) in each bracket}
\end{deriv}
\end{steps}

\answer{$P=Ke^{-r\tau}\Phi(\sigma\sqrt\tau-h)-S_0\Phi(-h)$, with the same $h$ as the call. This is \cardrow{card:bsprice}{2}.}

\section*{Part F: the copying portfolio at every node, and why it is self-financing}

Parts A and B found the \emph{price}. This part writes down the \emph{portfolio} that copies the claim, node by node, first in a two-period tree and then in a tree with any number of periods. In this part $\Delta$ is a number of shares (not the period length of Part D).

\begin{steps}[resume]
\item The bookkeeping. A portfolio is chosen at a node and then held, unchanged, for one period:
\begin{deriv}
\dl{\Delta=\text{number of shares bought at the node}}{kept until the next node}
\dl{b=\text{money put in the bond at the node}}{counted in money \emph{of the node's time}}
\dl{\text{cost at the node}=\Delta\,S+b}{$S$ is the stock price at the node; no $1+r$, you pay today's prices}
\dl{\text{worth one period later}=\Delta\,S_{\text{next}}+b(1+r)}{the shares are worth the new price; the bond money has grown for one period}
\end{deriv}
\begin{trap}
The $1+r$ goes on $b$ only, and only once for each period the money has sat in the bond. Never put $1+r$ on the shares: their growth is already inside the new price $S_{\text{next}}$. At the moment of buying there is no $1+r$ at all.
\end{trap}

\item The two-period tree and its names. Each node is named by the moves that lead to it:
\begin{deriv}
\dl{\text{time }0:\ S_0}{the root}
\dl{\text{time }1:\ uS_0\ (\text{node }u),\quad dS_0\ (\text{node }d)}{one move}
\dl{\text{time }2:\ u^2S_0,\quad udS_0,\quad d^2S_0}{two moves; up-down and down-up meet}
\dl{V_{uu}=g(u^2S_0),\ \ V_{ud}=g(udS_0),\ \ V_{dd}=g(d^2S_0)}{the claim's payoffs at time $2$}
\dl{(\Delta_0,b_0),\ (\Delta_u,b_u),\ (\Delta_d,b_d)}{the portfolio chosen at time $0$, at node $u$, at node $d$}
\end{deriv}
We do not know $V_u$, $V_d$, $V_0$ yet; they come out of the steps below. Work backwards, last period first.

\item The last period, at node $u$ (stock price $uS_0$). Choose $(\Delta_u,b_u)$ so that one period later it is worth the payoff in both cases:
\begin{deriv}
\dl{\Delta_u\cdot u^2S_0+b_u(1+r)=V_{uu}}{up again: shares worth $u\cdot uS_0$; bond grown once (Step~31)}
\dl{\Delta_u\cdot udS_0+b_u(1+r)=V_{ud}}{down: shares worth $d\cdot uS_0$}
\dl{\Delta_u=\frac{V_{uu}-V_{ud}}{u^2S_0-udS_0}}{subtract; $b_u(1+r)$ cancels}
\dl{b_u=\frac{V_{uu}-\Delta_u\,u^2S_0}{1+r}}{first line solved for $b_u$: divide by $1+r$, because $b_u$ is time-1 money that must \emph{grow into} what is needed at time 2}
\end{deriv}

\item The value at node $u$ is the cost of that portfolio at node $u$:
\begin{deriv}
\dl{V_u=\Delta_u\cdot uS_0+b_u}{Step~31: cost at the node, no $1+r$}
\dl{\phantom{V_u}=\frac{p^*V_{uu}+(1-p^*)V_{ud}}{1+r}}{the algebra of Steps 2--4 with $X_u=V_{uu}$, $X_d=V_{ud}$ and $S_0$ replaced by $uS_0$}
\end{deriv}

\item The same at node $d$ (stock price $dS_0$):
\begin{deriv}
\dl{\Delta_d\cdot udS_0+b_d(1+r)=V_{ud}}{up from node $d$: shares worth $u\cdot dS_0$}
\dl{\Delta_d\cdot d^2S_0+b_d(1+r)=V_{dd}}{down from node $d$: shares worth $d\cdot dS_0$}
\dl{\Delta_d=\frac{V_{ud}-V_{dd}}{udS_0-d^2S_0},\qquad b_d=\frac{V_{ud}-\Delta_d\,udS_0}{1+r}}{subtract; then the first line solved for $b_d$}
\dl{V_d=\Delta_d\cdot dS_0+b_d=\frac{p^*V_{ud}+(1-p^*)V_{dd}}{1+r}}{cost at node $d$; Steps 2--4 with $X_u=V_{ud}$, $X_d=V_{dd}$, $S_0$ replaced by $dS_0$}
\end{deriv}

\item The first period, at time $0$. The time-0 portfolio does \emph{not} have to pay the claim. It has to be worth exactly what the time-1 portfolio costs, $V_u$ at node $u$ and $V_d$ at node $d$:
\begin{deriv}
\dl{\Delta_0\cdot uS_0+b_0(1+r)=V_u}{at node $u$: worth one period later (Step~31) $=$ cost of $(\Delta_u,b_u)$}
\dl{\Delta_0\cdot dS_0+b_0(1+r)=V_d}{at node $d$: the same with $(\Delta_d,b_d)$}
\dl{\Delta_0=\frac{V_u-V_d}{uS_0-dS_0}}{subtract; $b_0(1+r)$ cancels}
\dl{b_0=\frac{V_u-\Delta_0\,uS_0}{1+r}}{first line solved for $b_0$}
\dl{V_0=\Delta_0S_0+b_0=\frac{p^*V_u+(1-p^*)V_d}{1+r}}{cost at time $0$; Steps 2--4 with $X_u=V_u$, $X_d=V_d$}
\end{deriv}

\item Self-financing at time 1. At node $u$ we sell the old portfolio and buy the new one. Compare the two at node $u$'s prices:
\begin{deriv}
\dl{\text{old, at node }u:\ \Delta_0\cdot uS_0+b_0(1+r)=V_u}{Step~31 (worth one period later); equals $V_u$ by the first line of Step~36}
\dl{\text{new, at node }u:\ \Delta_u\cdot uS_0+b_u=V_u}{Step~31 (cost at the node); equals $V_u$ by Step~34}
\dl{\text{old}-\text{new}=V_u-V_u=0}{no money is put in or taken out}
\dl{\text{old, at node }d:\ \Delta_0\cdot dS_0+b_0(1+r)=V_d}{second line of Step~36}
\dl{\text{new, at node }d:\ \Delta_d\cdot dS_0+b_d=V_d}{Step~35; again old $=$ new}
\end{deriv}
That is what self-financing means (\cref{card:selffin}): the old holding at the new prices equals the new holding at the new prices.

\item The whole strategy, in time order:
\begin{deriv}
\dl{\text{time }0:\ \text{pay }V_0=\Delta_0S_0+b_0}{the only money that ever goes in}
\dl{\text{time }1:\ \text{worth }V_u\text{ or }V_d}{Step~36}
\dl{\phantom{\text{time }1:\ }\text{switch to }(\Delta_u,b_u)\text{ or }(\Delta_d,b_d)\text{ for free}}{Step~37}
\dl{\text{time }2:\ \text{worth }V_{uu},\ V_{ud}\text{ or }V_{dd}}{the two lines of Step~33 or of Step~35}
\end{deriv}
So $V_0$ at the start turns into exactly the claim's payoff at the end, with nothing added on the way. Any other price for the claim is an arbitrage (Step~4's argument), so $V_0$ is the price.

\item Any number of periods. At time $t$ with stock price $S_t=s$, write $(\Delta_t[s],b_t[s])$ for the portfolio chosen there and $V_t[s]$ for the claim's value, with $V_T[s]=g(s)$. Steps 33--34 at this node:
\begin{deriv}
\dl{\Delta_t[s]\cdot us+b_t[s](1+r)=V_{t+1}[us]}{up: shares worth $us$; bond grown once}
\dl{\Delta_t[s]\cdot ds+b_t[s](1+r)=V_{t+1}[ds]}{down: shares worth $ds$}
\dl{\Delta_t[s]=\frac{V_{t+1}[us]-V_{t+1}[ds]}{us-ds}}{subtract}
\dl{b_t[s]=\frac{V_{t+1}[us]-\Delta_t[s]\cdot us}{1+r}}{first line solved for $b_t[s]$}
\dl{V_t[s]=\Delta_t[s]\cdot s+b_t[s]}{cost at the node (Step~31)}
\dl{\phantom{V_t[s]}=\frac{p^*V_{t+1}[us]+(1-p^*)V_{t+1}[ds]}{1+r}}{Steps 1--4 at price $s$; this is Step~6}
\end{deriv}

\item The two ways of writing $b$ are the same number:
\begin{deriv}
\dl{b_t[s]=V_t[s]-\Delta_t[s]\cdot s}{last line but one of Step~39: what is left of the cost after the shares; no $1+r$, all time-$t$ money}
\dl{b_t[s]=\frac{V_{t+1}[us]-\Delta_t[s]\cdot us}{1+r}}{Step~39: time-$(t+1)$ money brought back one period, so divide by $1+r$ once}
\end{deriv}
\cardrow{card:tree}{5} uses the first form. Use whichever you have the numbers for, and check the $1+r$ against the trap of Step~31.

\item Self-financing at every node. ``Old'' is the portfolio bought at $(t,s)$; ``new'' is the one bought at the next node. Let the stock go from $s$ at time $t$ to $us$ at time $t+1$ (the case $ds$ is the same with $d$ for $u$):
\begin{deriv}
\dl{\text{old: }\Delta_t[s]\cdot us+b_t[s](1+r)=V_{t+1}[us]}{first line of Step~39}
\dl{\text{new: }\Delta_{t+1}[us]\cdot us+b_{t+1}[us]=V_{t+1}[us]}{the cost line of Step~39 at the node $(t+1,us)$}
\dl{\text{old}=\text{new}}{no money in or out (\cardrow{card:tree}{6})}
\dl{\text{at }T:\ \Delta_{T-1}[s]\cdot S_T+b_{T-1}[s](1+r)}{the last portfolio, one period after it was bought}
\dl{\qquad=V_T[S_T]=g(S_T)}{first two lines of Step~39 with $t=T-1$: it pays the claim}
\end{deriv}
\end{steps}

\answer{At every node $(t,s)$ hold $\Delta_t[s]=\frac{V_{t+1}[us]-V_{t+1}[ds]}{us-ds}$ shares and $b_t[s]=V_t[s]-\Delta_t[s]\,s=\frac{V_{t+1}[us]-\Delta_t[s]\,us}{1+r}$ in the bond. It costs $V_t[s]$ when bought (no $1+r$) and is worth $V_{t+1}$ one period later (bond times $1+r$), which is exactly the cost of the next portfolio. So the strategy is self-financing, starts from $V_0$, and ends with the payoff $g(S_T)$.}
